% Auto-generated by ad-hoc script (see appendix on cath_share persistence).
% LaTeX preamble required: \usepackage{booktabs, threeparttable, siunitx}
\begin{table}[htbp]
\centering
\small
\begin{threeparttable}
\caption{Demographic stability of treatment groups, 1871--1890}
\label{tab:catholic_share_stability}
\begin{tabular}{l S S S S}
\toprule
 & {1871} & {1880} & {1885} & {1890} \\
\midrule
\multicolumn{5}{l}{\textit{Panel A. Prussia-wide Catholic share (pop-weighted)}} \\
Catholic share (\%)        & 33.564 & 33.744 & 33.977 & 34.227 \\
Counties observed ($N$)    & {514}  & {503}  & {547}  & {585}  \\
\midrule
\multicolumn{5}{l}{\textit{Panel B. Within-county means, balanced 4-census sample ($N=467$)}} \\
Low-Catholic ($\le 50\%$ in 1871, $N=314$)  & 10.463 & 10.824 & 11.005 & 11.330 \\
High-Catholic ($>50\%$ in 1871, $N=153$)    & 83.733 & 83.657 & 83.678 & 83.466 \\
\midrule
\multicolumn{5}{l}{\textit{Panel C. Within-county persistence (balanced clean sample, $N=467$)}} \\
$\rho\big(\text{cath\_share}_{1871},\, \text{cath\_share}_{1880}\big)$ & \multicolumn{4}{c}{0.99974} \\
$\rho\big(\text{cath\_share}_{1871},\, \text{cath\_share}_{1890}\big)$ & \multicolumn{4}{c}{0.99820} \\
$\rho\big(\text{cath\_share}_{1880},\, \text{cath\_share}_{1890}\big)$ & \multicolumn{4}{c}{0.99870} \\
\midrule
\multicolumn{5}{l}{\textit{Panel D. Sex ratio by treatment group (men per 100 women, pop-weighted)}} \\
Low-Catholic ($\le 50\%$ in 1871)   & 96.75 & 96.72 & 96.10 & 96.15 \\
High-Catholic ($>50\%$ in 1871)     & 97.36 & 97.25 & 96.69 & 96.76 \\
High $-$ Low (pp)                   & {+0.61} & {+0.53} & {+0.59} & {+0.61} \\
\bottomrule
\end{tabular}
\begin{tablenotes}
\footnotesize
\item \textit{Note:} Panel A reports Prussia-wide Catholic share computed as $\sum_i \text{Cath}_{i,t} / \sum_i \text{Pop}_{i,t}$, pooled across all counties present in the named census. Panels B and C restrict to the 475 counties observed in all four censuses; eight counties whose 1871--1890 change exceeds $\pm 17$pp are dropped as 1881 Prussian \textit{Kreis} boundary-reorganisation artefacts (verified to appear in administrative split/merge pairs), leaving $N=467$. Panel D uses all counties present in the named census ($N=348$--$494$ by year), with sex ratio constructed as $100 \times \sum_i \text{Pop}^M_{i,t} / \sum_i \text{Pop}^F_{i,t}$, treatment group fixed at the 1871 Catholic share. Sources: \texttt{REL1871}, \texttt{POP1871}, \texttt{POP1880}, \texttt{POP1885}, \texttt{POP1890} from Galloway's Prussian county-level files (\textit{Roman Catholic} only; the post-1870 Old Catholic schism is tabulated separately in 1890 as \texttt{RELDISSIDENT} and is not aggregated here). Definitions of ``Roman Catholic'' are stable across all four censuses; the 1880 sex ratio uses the de facto population (\texttt{Poplivingm/f}).
\item \textit{Reading the table.} Prussia-wide Catholic share drifts upward by 0.66pp across the 19-year window (Panel A). Within counties the drift is even smaller and goes in opposite directions for the two treatment groups: high-Catholic counties (Panel B, row 2) shift by $-$0.27pp (1871 $\to$ 1890), low-Catholic counties by $+$0.87pp. Persistence at the county level is essentially perfect: $\rho > 0.998$ over 19 years and $\rho > 0.999$ over the Kulturkampf-enforcement decade alone (Panel C). The 1871 $\to$ 1880 within-county standard deviation of cath\_share changes is just 0.87pp. Sex ratios (Panel D) are also stable: high-Catholic counties have a small but persistent surplus of men (+0.5 to +0.6 men per 100 women), and the high-low gap barely moves across the four censuses.
\item \textit{Identification implication.} The DiD treats $\text{cath\_share}_i$ as time-invariant from 1871 throughout 1862--1890. The data above confirm that assumption is empirically warranted: county-level Catholic shares do not respond meaningfully to the May Laws. Moreover, the timing rules out a conversion-response interpretation of the high-cath / low-cath fertility gap. If Catholics had converted under enforcement, we would expect the largest high-cath decline during 1871--1880; instead the (already small) decline is larger during the rollback period (1880--1890), consistent with secular industrialisation-driven migration rather than a policy-induced confessional response.
\item \textit{Sex ratio and the nuptiality gap.} Catholic counties have lower marriage rates and lower Coale $I_m$ than Protestant counties (e.g.\ $I_m^{1890} = 0.49$ vs $0.54$). Panel D shows that this cannot be a mechanical partner-availability effect: if anything, high-Catholic counties have \textit{more} men per woman, not fewer, and the 0.6pp gap is far too small (and the wrong sign) to mechanically explain the 4--5\% lower marriage rate. The Catholic nuptiality deficit is therefore a behavioural pattern (later age at marriage, higher share never-married, religious vocations) rather than a demographic-pool artefact, consistent with Hajnal's western European marriage pattern intensified by Catholic institutions.
\end{tablenotes}
\end{threeparttable}
\end{table}
